
在本章中，我们视情形或详或略地罗列本书所用的\trDifferential{}几何中的各类对象与工具。

\section{\trManifold{}与\trSubmanifold}
\label{sec:manifolds-submanifolds}

关于本节所要陈述的基本概念，可参见 [45]、[48]、[47]、[74]~与 [63]。

\subsection{\trManifold}
\label{subsec:manifolds}

一个\trManifold，就是一个在局部上极像 $\RR^{n}$ 的空间。让我们说得更加精确些。一个 $n$ 维\textbf{\trTopology\trManifold}（topological manifold）是一个分离的\trTopology{}空间，其每一点都含于一个\trHomeomorphism{}于 $\RR^{n}$ 的某个开子集的开子集之中。由这样的一个开集 $U$ 与一个\trHomeomorphism{} $\varphi$ 构成的偶对 $(U,\varphi)$，称为一个局部坐标系或\trChart。我们要考虑的\trManifold{}都至少是\trParacompact{}的，从而特别地是\trSeparable{}的。
\emph{我们}还要求 $V$ 是紧集的可数并，又是诸\trChart{}的开集的可数并。在本书中，我们要用的大多是\trCompact\trManifold，对它们而言这些性质自动成立。
一个\trManifold{}的诸\trChart{}所成之集称为一个\trAtlas。两个\trChart{} $(U,\varphi)$ 与 $(U',\varphi')$ 称为是\trCompatible{}的，如果复合
\[
\varphi(U\cap U')\xrightarrow{\ \varphi^{-1}\ }U\cap U'\xrightarrow{\ \varphi'\ }\varphi'(U\cap U')\text{，}
\]
——它把 $\RR^{n}$ 的一个开子集映成另一个开子集——是一个（$\cC^{\infty}$ 类的）\trDiffeomorphism。如果\trManifold{} $V$ 具有一个诸\trChart{}两两\trCompatible{}的\trAtlas，那么我们就称 $V$ 是一个\trDifferential\trManifold（或\trSmooth\trManifold，或 $\cC^{\infty}$ \trManifold）。

\noindent\textbf{例。}\trSphere{} $S^{n}$、诸\trProjectiveSpace{} $\CP^{n}(\RR)$ 与 $\CP^{n}(\CC)$，以及\trTorus{} $\RR^{n}/\ZZ^{n}$，都是（\trSmooth{}的）\trManifold。
按这个定义，一个零维\trManifold{}就是一个可数的离散集。

利用诸\trChart{}与 $\RR^{n}$ 的诸开子集上适当的概念，我们定义从一个\trManifold{}到另一个\trManifold{}的 $\cC^{\infty}$ 映射的概念：如果 $U$ 与 $V$ 分别含于\trManifold{} $M$ 与 $N$ 的诸\trChart{}的开子集之中，那么函数 $f:U\to V$ 是 $\cC^{k}$ 类的（对 $0\le k\le\infty$），如果复合
\[
\varphi(U)\xrightarrow{\ \varphi^{-1}\ }U\xrightarrow{\ f\ }V\xrightarrow{\ \psi\ }\psi(V)
\]
是 $\cC^{k}$ 类的。“$\cC^{k}$ 类”这一性质是局部的，而且由这个条件\trWellDefined。

\subsection{\trSubmanifold{}的刻画}
\label{subsec:submanifolds-characterization}

在诸\trManifold{}之中，有着 $\RR^{n}$ 的诸\trSubmanifold。$\RR^{n}$ 的一个\trSubmanifold{}是一个子空间 $V$，它在局部上极像一个线性子空间，其含义是：它的每一点在 $\RR^{n}$ 中都有一个开邻域 $U$，使之\trDiffeomorphic{}于 $\RR^{n}$ 中 $0$ 的一个邻域，这里所用的\trDiffeomorphism{} $\varphi$ 满足 $\varphi(U\cap V)=\varphi(U)\cap\RR^{d}$。整数 $d$ 就是\trSubmanifold{} $V$ 的维数。

\begin{theorem}\label{thm:submanifolds-characterization}
设 $V\subset\RR^{n}$。下列诸性质是等价的：
\begin{enumerate}[label=(\arabic*)]
\item $V$ 是 $\RR^{n}$ 的一个维数为 $d$ 的\trSubmanifold。
\item $V$ 的每一点 $x$ 在 $\RR^{n}$ 中都有一个开邻域 $U$，使得存在一个\trSubmersion{} $g:U\to\RR^{n-d}$ 满足 $U\cap V=g^{-1}(0)$。
\item $V$ 的每一点 $x$ 在 $\RR^{n}$ 中都有一个开邻域 $U$，使得存在 $\RR^{d}$ 中 $0$ 的一个邻域 $\Omega$ 与一个\trImmersion{} $h:\Omega\to\RR^{n}$，它是由 $\Omega$ 到 $U\cap V$ 上的一个\trHomeomorphism。
\end{enumerate}
\end{theorem}

换句话说，一个\trSubmanifold{}可以在局部上用方程来描述（此时方程的个数就是余维数），也可以用\trParametrization{}来描述（此时参数的个数就是维数）。所有这些性质都是\trImplicitFunctionTheorem{}与（或）\trLocalInversionTheorem{}的推论。
给定 $\RR^{n}$ 中、或者另一个\trManifold{}中一个被怀疑是\trSubmanifold{}的子集，要裁定这些判据中哪一个用来证明它的“罪名成立”最为顺手——这是一道出色且必不可少的练习。
\trManifold{}的\trSubmanifold{}的定义是类似的。
当然，\trSubmanifold{}都是\trManifold（这一点已包含在定义之中）。例如，一个实值函数的一个“\textbf{\trLevelSet}”（level set）——即某个实数 $\alpha$ 的\trPreimage——就落在性质~(2) 的适用范围之内：当 $f$ 沿 $f^{-1}(\alpha)$ 是一个\trSubmersion{}时，它是一个\trSubmanifold。
逆命题也成立：\trManifold{}也可以被视作\trSubmanifold，正如我们现在要回顾的那样。

\subsection{\trPartitionOfUnity、\trEmbedding{}定理}
\label{subsec:partitions-unity-embedding}

\emph{\trPartitionOfUnity。}
这是一个必不可少的工具，它使我们得以通过把注意力集中于诸结果的局部方面来证明许多结果。设 $V$ 是一个 $n$ 维\trCompact\trManifold\footnote{\trParacompact{}就足够了……}，并设 $(U_{i},\varphi_{i})_{1\le i\le N}$ 是 $V$ 的一个有限\trAtlas。我们可以证明：存在由\trOpen{}集 $\Omega_{i}$ 构成的 $V$ 的一个\trCovering，使得 $\overline{\Omega}_{i}\subset U_{i}$，还存在\trSupport{}位于 $U_{i}$ 内、且在 $\Omega_{i}$ 上取值为 $1$ 的函数 $h_{i}$。
令
\[
p_{i}=\frac{h_{i}}{\sum h_{i}}\text{，}
\]
我们便推出：对由有限多个\trOpen{}集 $\Omega_{1},\dotsc,\Omega_{N}$ 构成的 $V$ 的每一个\trCovering，都存在一族函数 $(p_{i})_{1\le i\le N}$，
\[
p_{i}:V\longrightarrow[0,+\infty[
\]
使得
\[
\operatorname{Supp}(p_{i})\subset\Omega_{i}\quad\text{且}\quad\sum_{i=1}^{N}p_{i}\equiv 1\text{。}
\]

\emph{\trCompact\trManifold{}在\trEuclidean{}空间中的\trEmbedding。}
设 $V$ 是一个 $n$ 维\trCompact\trManifold，并设 \allowbreak$(U_{i},\varphi_{i})_{1\le i\le N}$ 是 $V$ 的一个有限\trAtlas。我们利用上面引入的诸函数 $h_{i}$ 来构造一个映射 $F$，其定义为
\[
F=(h_{1}\varphi_{1},\dotsc,h_{N}\varphi_{N},h_{1},\dotsc,h_{N}):V\longrightarrow(\RR^{n})^{N}\times\RR^{N}=\RR^{(n+1)N}\text{。}
\]
我们可以验证：$F$ 是一个\trInjective\trImmersion，因而是一个\trEmbedding；参见下述定理。

\begin{theorem}[\trWhitney{}\trEmbedding{}定理]\label{thm:whitney-embedding-theorem}\BookIndexDefinition{T013}
设 $V$ 是一个\trCompact\trManifold；那么存在把 $V$ 作为\trSubmanifold{}\trEmbedding{}进一个\trEuclidean{}空间 $\RR^{N}$ 的\trEmbedding。
\end{theorem}
\noqed

利用一个与我们稍后要回顾的论证相类似的一般位置论证，我们可以证明：\trEuclidean{}空间的维数可以约减到 $2\dim V+1$。

\noindent\textbf{注记。}这个结果属于实\trDifferential{}几何；它在复解析几何中没有类似物：由\trLiouville{}定理，任何\trCompact{}的解析\trManifold{}都不是 $\CC^{N}$ 的解析\trSubmanifold。

\subsection{边界}
\label{subsec:boundaries}

一个\trManifoldWithBoundary{}是一个\trTopology{}空间 $V$，其每一点都有一个\trHomeomorphism{}于 $\RR^{n}$ 中半空间 $(x_{n}\le 0)$ 的某个开子集的邻域。那些不具有\trHomeomorphism{}于 $\RR^{n}$ 的一个开子集的邻域的点，就是边界点。边界记作 $\partial V$，是一个 $n-1$ 维\trSubmanifold。我们可以证明（参见\cref{sec:vector-fields-odes}），边界具有形如 $\partial V\times[0,\varepsilon[$ 的邻域。

\subsection{切向量、\trTangentMap}
\label{subsec:tangent-vectors-tangent-maps}

让我们回忆一下：我们可以把\trManifold{} $V$ 在 $x$ 处的一个切向量定义为 $V$ 上过 $x$ 点的诸曲线所成的一个等价类，即
\[
c:\ ]-\varepsilon,\varepsilon[\ \longrightarrow V\quad\text{满足 }c(0)=x\text{，}
\]
所用的等价关系是：如果对于一张以 $x$ 为中心的\trChart{} $\varphi$，我们有
\[
(\varphi\circ c_{1})'(0)=(\varphi\circ c_{2})'(0)\text{，}
\]
就认同 $c_{1}$ 与 $c_{2}$。
在 $x$ 处的诸切向量所成之集是一个\trVectorSpace，我们把它记作 $T_{x}V$。这使我们得以定义\trTangentMap
\[
T_{x}f:T_{x}V\longrightarrow T_{f(x)}W
\]
它与映射 $f:V\to W$ 相伴：做法无非是把曲线 $c$ 的类送到曲线 $f\circ c$ 的类。
它也使我们得以把一个切向量视为函数上的一个\trDerivation，也就是说，视为一个线性映射 $X$，它定义在 $x$ 的邻域内的 $\cC^{\infty}$ 函数之集上，且满足
\[
X(fg)=f(x)X(g)+g(x)X(f)\text{。}
\]
对每个 $x$，给定一个向量 $X\in T_{x}V$ 与一个函数 $f:V\to\RR$，我们令 $X(f)(x)=T_{x}f(X)\in\RR$；这就定义了函数 $X(f)$。

\begin{remark}\label{thm:df-notation}
这个函数常常记作 $(df)_{x}$ 而非 $T_{x}f$，尤其当 $f$ 是一个实值函数时如此。在此我们就将这样做。
\end{remark}

用每一种刻画来描述一个\trSubmanifold{}的\trTangentSpace，都是相当有教益的（特别地，我们会发现：当 $f$ 是\trSubmersion{}时，$T_{x}\bigl(f^{-1}(0)\bigr)=\operatorname{Ker}T_{x}f$）。

\subsection{\trVectorField}
\label{subsec:vector-fields}

我们还定义了\textbf{\trTangentBundle}（tangent bundle）\footnote{同样地，利用诸\trTangentSpace{}的各个对偶，也可以定义\trCotangent{}丛。}，做法是把\trManifold{}在各点处的诸\trTangentSpace{}的\trDisjointUnion{}赋予一个\trTopology{}与适当的\trManifold{}结构。
诸\trVectorField{}是\trTangentBundle{}的 $\cC^{\infty}$ \trSection\footnote{一个\trVectorBundle{} $\pi:E\to V$ 的\trSection{}是一个映射 $s:V\to E$，满足 $\pi\circ s=\mathrm{Id}_{V}$。}。换句话说，一个\trVectorField{}就是：对每一点，给出一个以 $\cC^{\infty}$ 方式依赖于所论点的切向量。一般地，我们将用 $X$ 这类罗马大写字母来表示\trVectorField，而 $X_{x}$ 将表示场 $X$ 在点 $x$ 处的取值。
如果 $X$ 是一个\trVectorField{}而 $f$ 是一个函数，那么我们用 $X\cdot f$ 表示 $f$ 沿 $X$ 的导数，其定义为
\[
(X\cdot f)(x)=X_{x}\cdot f=T_{x}f(X_{x})\text{。}
\]
一切\trDerivation{}都由\trVectorField{}产生。

\emph{\trLieBracket。}
两个\trDerivation{}的复合不再是一个\trDerivation；然而，
\[
f\longmapsto X\cdot(Y\cdot f)-Y\cdot(X\cdot f)
\]
却是一个。我们把与之相应的\trVectorField{}记作 $[X,Y]$。
一个\trVectorField{}也可以被视作\trManifold{}上的一个\trDifferential{}方程。参见\cref{sec:vector-fields-odes}。

\subsection{\trDifferential\trForm}
\label{subsec:differential-forms}

\trManifold{} $V$ 上的一个 $p$ 次\trDifferential\trForm{}（也称为一个 $p$-形式）是丛 $\Lambda^{p}T^{\star}V$ 的一个\trSection。用不那么做作的话来说，它就是：对每个 $x\in V$，给出 $T_{x}V$ 上的一个 $p$ 次外\trDifferential\footnote{译注：原书此处印作“an exterior derivative of degree $p$ on the tangent space $T_{x}V$”；按上下文——$p$-形式是丛 $\Lambda^{p}T^{\star}V$ 的\trSection——此处所指应是 $T_{x}V$ 上的一个交替 $p$-线性形式（即一个 $p$ 次外形式），并非把 $p$-形式变成 $(p+1)$-形式的外\trDifferential{}算子 $d$（见本节下文），疑为用词之疏。此处照原书排印，存疑备案。——译者注}，它以 $\cC^{\infty}$ 方式依赖于 $x$。在 $V$ 的局部坐标 $(x_{1},\dotsc,x_{n})$ 下，它可以写成
\[
\alpha=\sum\alpha_{i_{1},\dotsc,i_{p}}\,dx_{i_{1}}\wedge\dotsb\wedge dx_{i_{p}}\text{，}
\]
其中诸 $\alpha_{i_{1},\dotsc,i_{p}}$ 是函数，而 $dx_{1},\dotsc,dx_{n}$ 表示 $(\RR^{n})^{\star}$ 的\trCanonical{}基。
一个映射 $f:V\to W$ 通过
\[
(f^{\star}\alpha)_{x}(X_{1},\dotsc,X_{p})=\alpha_{f(x)}\bigl(T_{x}f(X_{1}),\dotsc,T_{x}f(X_{p})\bigr)
\]
把 $W$ 上的诸\trDifferential\trForm{}送到 $V$ 上的诸\trDifferential\trForm。
作用于诸\trDifferential\trForm{}的最重要的运算是外\trDifferential。它把一个 $p$-形式变成一个 $(p+1)$-形式；在坐标下，我们有
\[
d\alpha=\sum d\alpha_{i_{1},\dotsc,i_{p}}\wedge dx_{i_{1}}\wedge\dotsb\wedge dx_{i_{p}}
\]
（其中 $d\alpha_{i_{1},\dotsc,i_{p}}$ 是函数 $\alpha_{i_{1},\dotsc,i_{p}}$ 的\trDifferential）。
它满足 $d\circ d=0$。有\trClosed\trForm（$d\alpha=0$）与\trExact\trForm（$\alpha=d\beta$）。\trExact\trForm{}都是\trClosed{}的，但逆命题一般不成立——这就产生了\trRham\trCohomology。
我们特别需要的是 $1$-形式与 $2$-形式、外\trDifferential、\trClosed\trForm、\trExact\trForm、\trVolume\trForm{}的概念以及\trStokes{}定理。参见 [45,~Chapter~V]~与 [74,~Chapter~VII]。

\subsection{\trOrientation{}事宜}
\label{subsec:orientation-matters}

如果\trManifold{}具有一个有向\trAtlas（即其所有转移映射都保持\trOrientation），我们就称它是\trOrientable{}的。具有一个\trVolume\trForm{}的\trManifold{}是\trOrientable{}的（证明是立时的）。逆命题也成立：我们可以利用一个\trPartitionOfUnity{}来构造一个\trVolume\trForm（\trManifold{}为\trCompact{}时参见 [45,~Chapter~V]，一般情形参见 [37,~p.~46]）。
有向\trManifold{}的边界具有一个自然的\trOrientation。
\noindent\textbf{例。}如果 $f:\RR^{n}\to\RR$ 是一个\trSubmersion，那么\trSubmanifold{} $f^{-1}(0)$ 是\trOrientable{}的。为了看出这一点，我们可以利用该函数的\trGradient——它是一个垂直于该\trSubmanifold{}的\trVectorField。
